\paragraph{How the hypotheses are made simultaneous.}\label{assembly:selection}
The quantifier on the nodes matters here. Theorem~A chooses the weights
at scale $x$ before $u_0$ and requires a family of nodes, one for each
$u_0$. A bound for independently prescribed nodes would be stronger.
The original assembly constructs the nodes and their arithmetic
identifications together.

Fix the primitive character. A family $\mathcal F$ has a window and,
for each $x$, an index set of records. The construction uses finite
collections of lower bounds on parameters, some depending on earlier
parameters. Call such a collection $K$. Two collections have a
joint bound $K\vee K'$ by taking maxima at each prescribed argument.
A family satisfying the joint bound satisfies each separately, with
the same height parameter $C_H$. A family assertion $P$ used in the
construction has the following quantifiers. There is an $\epsilon_P>0$ such that for each $0<\varepsilon\le\epsilon_P$
there is $K_P$ for which every family in the specified class, large for
$(K_P,C_H)$, satisfies $P(\mathcal F,C_H)$.
Two assertions of this form hold jointly by taking the smaller
$\epsilon_P$ and the joint parameter bound. Repeat for the following
nine assertions: the local comparison, the conditional division
estimate, the $\eta$ estimate, irreducibility, equality of ranks at
small level, fullness and the slope-rank identity, the gain estimate,
the slope estimate, and the global accounting implication.
The design-existence theorem, for each sufficiently small
$\varepsilon$ and \emph{every} joint bound, supplies one family and
one $C_H$ meeting that bound, the selection condition and the two size
conditions on the auxiliary primes and the determinant. Apply it
after forming the joint bound. All nine implications can now be
evaluated at precisely this family and $C_H$.

On this family the window estimate, local comparison and small-level
rank equality give divisibility at the twisted places. Insert this and
the $\eta$ estimate into the conditional division estimate. The
resulting division data together with fullness, the rank identity
and the two size conditions give the required slope comparison.
The global implication consumes the gain, slope, irreducibility,
division and slope-comparison conclusions. In the original application
the parameters for this join are $C_{\rm abs}=2$, $C_0=3$, $e_D=2$
and $e_L=0$. The family and its bounds are fixed before the remaining
quantifiers over $x$. The record exists at every sufficiently large
$x$ in that family's window.

At such an $x$, one index $i_0$ fixes the frame, $m,f$, and the host
primes. One function $u_0\mapsto i(u_0)$ selects the other indices.
At $i(u_0)$ the frame, ranks and two host primes agree with those at
$i_0$, the marked element is $u_0$, and two integer functions supply
the top weights at every index. The global identities fix, among
other entries, $v$, the logarithmic derivative, and
$g=\log(g_0g_1)$; this is the one $g$ in (A5) and (A6).
The local identities read off the ranks, primes, lattice slopes,
heights and ramification-weighted index losses at $i(u_0)$.
The arithmetic identities identify the top and thick index sets and say
that whenever an embedding has the code of a node, that node equals
the archimedean parameter of this very index and lift. The package
also includes the size bounds, error estimates, and identifications
of the finite indexing sets for the sums. Consequently the data in
the estimates and in the final excess sum agree at every occurrence.

The source stores eleven conditions in the size record at this point:
the Hida setting, division bound, summed local error, height comparison,
word count, slopes, quotient-slope bounds, losses, slope inequality,
height bound and parameter bounds. The slope inequality is evaluated at
the one height stored on the record; its upper bound is a separate field.
The local-error entry bounds $|E_u-mAg|$. This suffices both when
comparing two words and when using (A5).

For the last comparison, use the chosen bijections from the top labels
to $T$ and from the archimedean indices to $\Phi$. Reindexing both
finite sums makes the arithmetic excess exactly the accounted sum
minus its main term. The size and error bounds on this record give
$x^{(e_D+e_L+3)\varepsilon}\mathcal B$. The window prime-sum estimate
costs $3\varepsilon$ more in the exponent and absorbing the fixed
coefficients costs a further $\varepsilon$. Thus the one-sided bound
has exponent $1/2+(e_D+e_L+7)\varepsilon=1/2+9\varepsilon$
on the selected windows. For each $s$ with $\operatorname{Re}s>1/2$,
choose $\varepsilon$ small enough that this exponent is below
$\operatorname{Re}s$. Corollary~\ref{corA} then gives nonvanishing.
This order of choices requires arbitrarily large windows of the selected
family and uniformity in $x$ throughout each such window.

For the manuscript formulation it remains to verify that the lattices,
period maps, metrics, nodes and local lengths constructed in the
preceding sections satisfy these identifications together. For example,
Paragraph~\ref{assembly:wordheight} computes $h_u$ on the selected record. Its comparison with the independently defined periods must give the
needed inequalities for the quantity
in Definition~\ref{r4:R4:hu}; that comparison enters both (K4) and
(K6). Proposition~\ref{a8:prop:record} then supplies the fixed-record
inequalities, and Subsection~\ref{r6:R6:sec:thmA} and
Theorem~\ref{mainthm} give the stated implication to nonvanishing.
For an imprimitive character the missing Euler factors are nonzero
when $\operatorname{Re}s>0$. Thus the primitive result gives the
all-character nonvanishing input, away from the pole for the principal
character. For zeta, a nontrivial zero to the left of $1/2$ would give
one to its right by the functional equation. This gives the terminal
Riemann hypothesis assertion.

\paragraph{The height at one index.}\label{assembly:wordheight}
Fix an index of the selected family. Its height identity states, for
each word $u$ and for the same constant $C=C_{\rm sz}$, that its
stored height is
\[
 h^{\rm rec}_u=P_u+\sum_{j=1}^m K_{j,u}
             -\sum_{j=1}^m L_{6,j}
             -[K:\mathbb Q]^{-1}\sum_{\sigma:K\hookrightarrow\mathbb C}
                  \log\kappa_\sigma .
\]
Here $P_u=-A\log T_u+e_P^{\max}$ with the record's $C$ in
$e_P^{\max}$. Write $N_j=k+j(t_r+1)$, $N=N_m$, and denote the
record's row weights by $\nu_{j\rho i}(u)$. Then
\begin{align*}
 K_{j,u}
   &=\sum_\rho\sum_{0\le i<N_{j-1}}
       \bigl(\ph(\nu_{j\rho i}(u))-\nu_{j\rho i}(u)\log\Omega\bigr)
       +CfN^2\Lambda_a+Cfm^3N^2L_\varepsilon,\\
 L_{6,j}
   &=\sum_\rho\log H^+_{j\rho}-Cfn_V^2 L_p(b).
\end{align*}
The last expression uses the hook-length product for the positive
partition of Paragraph~\ref{frame:forms}. In factorial coordinates its
logarithm for rows $a_0,\ldots,a_{n-1}$ is
$\sum_{i<n}\log\Gamma(a_i+n-i)-
\sum_{i<l<n}\log(a_i-a_l+l-i)$. The two expressions agree by the
hook-length formula. The mean over the embeddings agrees with the
degree-weighted mean over infinite places when
$\kappa_\sigma=\kappa_{v(\sigma)}$. Substitution proves the height
identity with these inputs at this index.

In $(h^{\rm rec}_u+A\log T_u)-(h^{\rm rec}_{u'}+A\log T_{u'})$
the pairing error, the two uniform kernel errors, the hook terms and
the mean of $\log\kappa$ cancel. In the remaining double sum only
the rows $1\le i<k-1$ can depend on the word. At these rows the two
weights are at least one, differ by at most the small-weight bound
$S$, and are at most $M=6x^6b$ on the selected scale. On $[1,M]$ the
derivative of $\ph$ is $\log y$, bounded by $\log M$. The mean-value
inequality therefore bounds the difference of one summand by
$S(\log M+|\log\Omega|)$. Summing gives
$mfkS(\log M+|\log\Omega|)$. The size estimates on this family give
$kS\le b$ and
$\log M+|\log\Omega|\le(14+6/\varepsilon)L_\varepsilon$.
Consequently the variation is at most
$(20+6/\varepsilon)(N_Ta\Lambda_a+mB_0)$, since
$B_0=fmbL_\varepsilon$ and $m\ge1$. The constant is independent of
the word and scale. This proves (K6) for the selected height and
explains the word index suppressed in $s_j$ when only one word is fixed.

For Definition~\ref{r4:R4:hu} the corresponding term-by-term test is
now explicit. Evaluate its row sums on this realization, padding the
$N_{j-1}$ rows with zeros if needed, use its positive hook factors,
and match its mean of $\log\kappa$. Replacing the sharp errors by
larger errors independent of $u$ adds the same constant to every
$h_u$. It cancels from (K6). In (K4) it adds exactly that constant
to the left-hand side, so its bound must also be included in the
right-hand error. For the frame term this absorption can be made
explicit. Keep the other inputs fixed and replace the per-lift
$Cfn_V^2L_p(b)$ by $C_\lambda fN^2\Lambda_a$ with
$C_\lambda\ge20C$. Paragraph~\ref{frame:forms} makes this a
majorant. Its total increment is at most
$C_\lambda mfN^2\Lambda_a\le 2C_\lambda mfbL_\varepsilon
=2C_\lambda B_0$, using $N^2\le b$ and
Lemma~\ref{r6:R6:lem:sizes}(iii). This is absorbed in the
$m(B_0+N_c)x^{e_D\varepsilon}$ error of (K4) by increasing its
constant by $2C_\lambda$. The height of the independently defined period map
still requires the metric and pairing identifications of
Paragraph~\ref{fr:sec:heights}.

